\documentclass[a4paper,12pt]{article}
\usepackage[utf8]{inputenc}
\usepackage[italian]{babel}
\usepackage{listings}
\usepackage{xcolor}
\usepackage{float}
\usepackage{placeins}
\usepackage{url}
\setlength{\emergencystretch}{3em}

\lstset{
  language=C,
  backgroundcolor=\color{lightgray!20},
  frame=single,
  breaklines=true,
  numbers=left,
  numberstyle=\tiny,
  stepnumber=1,
  numbersep=5pt,
  captionpos=b,
  basicstyle=\ttfamily\small,
  keywordstyle=\color{blue}\bfseries,
  commentstyle=\color{green!60!black},
  stringstyle=\color{red},
  showstringspaces=false
}

\title{LabReSiD26 -- Hands-On 12}
\author{Davide Ferrara -- 518629}
\date{}

\begin{document}

\maketitle

% ---------------------------------------------------------------------------
\section{Esercizio 1: La shell base}
% ---------------------------------------------------------------------------

In questo esercizio si implemeenta un ciclo \emph{read-fork-exec-wait}:
l'obbiettivo è leggere un comando da \texttt{stdin},
fare \texttt{fork()}, eseguire il programma nel figlio con
\texttt{execvp()} e aspettarne la terminazione nel
padre con \texttt{waitpid()}.

\subsection*{Prompt e lettura}

Il prompt viene stampato con \texttt{printf} seguito da \texttt{fflush(stdout)}:
\texttt{printf} usa un buffer \emph{line-buffered} e, poiché \texttt{"minish> "}
non termina con \texttt{'\textbackslash n'}, senza \texttt{fflush} il testo
resterebbe nel buffer e non apparirebbe sullo schermo prima che il programma
si blocchi su \texttt{fgets()}.

\subsection*{Figlio: \texttt{execvp} e padre: \texttt{waitpid}}

\texttt{execvp(argv[0], argv)} sostituisce il processo con il
programma richiesto cercandolo nel \texttt{PATH}. Se ha successo non ritorna
mai: il codice dopo è raggiungibile solo in caso di errore, dove si termina
con \texttt{exit(127)} (convenzione POSIX per \emph{command not found}).
Il padre chiama \texttt{waitpid(pid, \&status, 0)} e decodifica il risultato
con le macro \texttt{WIFEXITED}/\texttt{WEXITSTATUS} e
\texttt{WIFSIGNALED}/\texttt{WTERMSIG}.

\begin{lstlisting}[caption={Figlio ed padre -- parti cruciali}]
if (pid == 0) {
    /* execvp non ritorna mai in caso di successo */
    if (execvp(argv[0], argv) == -1) {
        perror(argv[0]);
        exit(127);
    }
} else {
    /* waitpid() aspetta il figlio specificato */
    int status;
    waitpid(pid, &status, 0);
    if (WIFEXITED(status) && WEXITSTATUS(status))
        printf("[Exit %d]\n", WEXITSTATUS(status));
    if (WIFSIGNALED(status))
        printf("[Signal %d]\n", WTERMSIG(status));
}
\end{lstlisting}

\subsection{Output atteso}
Come si può notare dall'output, la shell è in grado di eseguire \textbf{ls}, di fare
\textbf{echo} e di uscire con \textbf{exit code 127}.

\begin{lstlisting}[language=bash, caption={Test Esercizio 1}]
$ ./minish_es1
minish> ls -l
total 472
-rwxr-xr-x 1 dff dff  15784 May 13 crash
-rw-r--r-- 1 dff dff     63 May 13 crash.c
-rwxr-xr-x 1 dff dff  16600 May 13 minish_es1
minish> echo ciao mondo
ciao mondo
minish> comando_inesistente
comando_inesistente: No such file or directory
[Exit 127]
minish> exit
\end{lstlisting}

% ---------------------------------------------------------------------------
\section{Esercizio 2: Zombie, segnali e \texttt{cd}}
% ---------------------------------------------------------------------------

L'esercizio estende la shell con tre funzionalità: esecuzione in background
con \texttt{\&}, raccolta asincrona dei figli tramite \texttt{SIGCHLD} e il
built-in \texttt{cd}.

\subsection*{Comandi in background con \texttt{\&}}

Se l'ultimo token della riga è \texttt{\&}, si rimuove da \texttt{argv},
si stampa il PID del figlio e si torna al loop senza chiamare
\texttt{waitpid()}.

\begin{lstlisting}[caption={Rilevazione background e handler SIGCHLD}]
/* Il loop e' necessario: i segnali dello stesso tipo non si accodano,
 * un singolo SIGCHLD puo' rappresentare N terminazioni. */
static void sigchld_handler(int sig) {
    (void)sig;
    while (waitpid(-1, NULL, WNOHANG) > 0)
        ;
}

...

/* In main, prima del loop: */
signal(SIGCHLD, &sigchld_handler);

...

/* Rilevazione &: */
int bg = argc > 0 && !strcmp(argv[argc - 1], "&");
if (bg)
    argv[--argc] = NULL;

/* Nel padre: */
if (bg) { printf("[1] %d\n", pid); continue; }
\end{lstlisting}

\subsection*{Raccolta asincrona con \texttt{SIGCHLD}}

Quando un figlio in background termina, il kernel manda \texttt{SIGCHLD} al
padre. L'handler chiama \texttt{waitpid(-1, NULL, WNOHANG)} in un loop per
raccogliere tutti i figli terminati, evitando i processi zombie.

\medskip
\noindent\textbf{Senza handler} -- zombie accumulati:
\begin{lstlisting}[language=bash]
dff   73704  S+  ./minish
dff   73731  Z+  [sleep] <defunct>
dff   73783  Z+  [sleep] <defunct>
...  (15 zombie)
\end{lstlisting}

\noindent\textbf{Con handler} -- nessun zombie:
\begin{lstlisting}[language=bash]
dff   72437  S+  ./minish
dff   73119  S+  grep -E minish|defunct
\end{lstlisting}

\subsection*{Built-in \texttt{cd}}

\texttt{chdir()} modifica la directory corrente del processo che la chiama.
Se lo facessimo in un figlio, solo la directory del figlio cambierebbe; la
shell rimarrebbe nella directory di partenza. Per questo \texttt{cd} viene
gestito prima della \texttt{fork()}, con \texttt{continue} per saltarla.
Senza argomenti va in \texttt{\$HOME}.

\begin{lstlisting}[caption={Built-in cd}]
if (!strcmp(argv[0], "cd")) {
    if (argc == 1) {
        const char *home = getenv("HOME");
        if (!home) { fprintf(stderr, "cd: HOME non impostata\n"); continue; }
        if (chdir(home) < 0) perror("cd");
    } else {
        if (chdir(argv[1]) < 0) perror("cd");
    }
    continue; /* non fare fork */
}
\end{lstlisting}

% ---------------------------------------------------------------------------
\section{Esercizio 3: La pipe}
% ---------------------------------------------------------------------------

L'esercizio estende la shell per supportare \texttt{cmd1 | cmd2}: i due
comandi vengono eseguiti in due processi figli collegati da una pipe del kernel.

\subsection*{Rilevazione di \texttt{|} e split di \texttt{argv[]}}

Si scorre \texttt{argv[]} finché non si trova \texttt{"|"}: si imposta
\texttt{argv[i] = NULL} per terminare il vettore sinistro e si salva
\texttt{argv\_right = \&argv[i+1]} come puntatore al comando destro.

\begin{lstlisting}[caption={Ricerca di | e pipeline}]
char **argv_right = NULL;
int i = 0;
while (i < argc) {
    if (!strcmp(argv[i], "|")) {
        argv[i] = NULL; // Split
        argv_right = &argv[i + 1]; // Argomenti di destra
        break;
    }
    i++;
}

\end{lstlisting}

\subsection*{Esecuzione della pipeline}

Prendendo come esempio \texttt{ls | grep .c}: \texttt{grep} deve leggere
da \texttt{stdin} l'output di \texttt{ls}. Si creano due processi figli:
il primo redirige \texttt{stdout} sul capo di scrittura della pipe
(\texttt{fd[1]}), il secondo redirige \texttt{stdin} sul capo di lettura
(\texttt{fd[0]}). Gli indici della pipe rispecchiano la convenzione dei
file descriptor standard: \texttt{fd[1]} è il capo di scrittura (come
\texttt{stdout}\,=\,1) e \texttt{fd[0]} è il capo di lettura (come
\texttt{stdin}\,=\,0).

\begin{lstlisting}[caption={Fork, redirect e attesa della pipeline}]
pid_t child1 = fork();
if (child1 == 0) {
    dup2(fd[1], STDOUT_FILENO); /* stdout -> pipe */
    close(fd[0]); close(fd[1]);
    execvp(argv_left[0], argv_left);
    exit(127);
}
pid_t child2 = fork();
if (child2 == 0) {
    dup2(fd[0], STDIN_FILENO);  /* stdin <- pipe */
    close(fd[1]); close(fd[0]);
    execvp(argv_right[0], argv_right);
    exit(127);
}
/* Il padre chiude entrambi i capi: se non chiudesse fd[1],
 * child2 non riceverebbe mai EOF e rimarrebbe bloccato. */
close(fd[0]); close(fd[1]);
waitpid(child1, NULL, 0);
waitpid(child2, NULL, 0);
\end{lstlisting}

\newpage
\subsection{Output atteso}

\begin{lstlisting}[language=bash, caption={Test Esercizio 3 -- pipe}]
minish> ls -1 | grep .c
crash.c
minish.c
minish_es1.c
minish_es2.c
minish> cat /etc/passwd | grep root
root:x:0:0::/root:/usr/bin/bash
minish> echo uno due tre | wc -w
3
minish> ps aux | grep minish
dff  110138  0.0  0.0  2752  1792 pts/3  S+  23:18  0:00 ./minish
dff  111198  0.0  0.0  6676  4212 pts/3  S+  23:24  0:00 grep minish
\end{lstlisting}

\end{document}
